% Source-derived chapter 02: Perverse filtrations
% Original source: hitchin-fibrations-abelian-surfaces-pw
\chapter{Perverse filtrations}
\label{hitchin-fibrations-abelian-surfaces-pw-chapter-02}
\source{hitchin-fibrations-abelian-surfaces-pw}

\begin{remark}[Source section]
\label{hitchin-fibrations-abelian-surfaces-pw-chapter-02-source}
\source{hitchin-fibrations-abelian-surfaces-pw}
\source{hitchin-fibrations-abelian-surfaces-pw}
This chapter reproduces the corresponding section of the retrieved
LaTeX source, including its definitions, statements, examples, and
proofs. It has not yet been translated into Lean.
\notready
\end{remark}

% The source section title was: Perverse filtrations
\label{sec1}
\source{hitchin-fibrations-abelian-surfaces-pw}

\subsection{Overview}
We begin with the definition and some relevant properties of the perverse filtration associated with a proper surjective morphism $\pi: X\rightarrow Y$. Some  references are \cite{BBD, dCM0, dCM1, WSB}. 
Throughout this section, for simplicity, we assume $X$ and $Y$ are nonsingular.

Following \cite{dCHM3, Z, SZ}, we discuss the perverse filtration for the Hilbert scheme of $n$-points on an abelian surface $E \times E'$ product of two elliptic curves, induced by the natural  second projection morphism
\[
E \times E' \to E'.
\]
 This example plays a crucial role in Section 2 concerning moduli of sheaves on abelian surfaces.

\subsection{Perverse sheaves}
Let $D_c^b(Y)$ denote the bounded derived category of $\BQ$-constructible sheaves on $Y$, and let $\BD: D_c^b(Y)^{\mathrm{op}} \to D_c^b(Y)$ be the Verdier duality functor. The full subcategories 
\begin{align*}
    ^ \mathfrak{p} D_{\leq 0}^b(Y) & = \left\{\CE \in D_c^b(Y): \dim \mathrm{Supp}(\CH^i(\CE)) \leq -i \right\}, \\
    ^ \mathfrak{p} D_{\geq 0}^b(Y) & = \left\{\CE \in D_c^b(Y): \dim \mathrm{Supp}(\CH^i(\BD\CE)) \leq -i \right\}
\end{align*}
give rise to the \emph{perverse $t$-structure} on $D_c^b(Y)$, whose heart is the abelian category of perverse sheaves,
\[
\mathrm{Perv}(Y) \subset D_c^b(Y).
\]

For $k \in \BZ$, let $^\mathfrak{p}\tau_{\leq k}$ be the truncation functor associated with the perverse $t$-structure. Given an object $\CC \in D_c^b(Y)$, there is a natural morphism
\begin{equation}\label{trun0}
\source{hitchin-fibrations-abelian-surfaces-pw}
^\mathfrak{p}\tau_{\le k}\CC \rightarrow \CC.
\end{equation}
For the morphism $\pi:X \to Y$, we obtain from (\ref{trun0}) the morphism
\[
^\mathfrak{p}\tau_{\le k}R\pi_\ast \BQ_X \rightarrow R\pi_\ast \BQ_X,
\]
which further induces a morphism of (hyper-)cohomology groups
\begin{equation}\label{perv_filt}
\source{hitchin-fibrations-abelian-surfaces-pw}
H^{d-(\dim X -{ R})}\Big{(}Y, \, ^\mathfrak{p}\tau_{\le k}(R\pi_\ast \BQ_X[\dim X - { R}]) \Big{)} \rightarrow H^d(X, \BQ).
\end{equation}
Here
\begin{equation}\label{defect}
\source{hitchin-fibrations-abelian-surfaces-pw}
    {R} = \dim X \times_Y X - \dim X
\end{equation}
is the \emph{defect of semismallness}. The $k$-th piece of the perverse filtration (\ref{Perv_Filtration})
\[
P_kH^d(X, \BQ) \subset H^d(X, \BQ)
\]
is defined to be the image of (\ref{perv_filt}).\footnote{Here the shift $[\dim X- {R}]$ is to ensure that the perverse filtration is concentrated in the degrees $[0, 2 { R}]$.} 

We say that a  graded vector space decomposition
\[
H^*(X, \BQ) = \bigoplus_{k,d} G_k H^d(X, \BQ)
\]
\emph{splits} the perverse filtration, if 
\[
P_k H^d(X, \BQ) = \bigoplus_{i \leq k} G_i H^d(X, \BQ), \quad \forall k,d.
\]


\subsection{The decomposition theorem}\label{sec1.3}
\source{hitchin-fibrations-abelian-surfaces-pw}
Perverse filtrations can be described through the decomposition theorem \cite{BBD, dCM0}. By applying the decomposition theorem to the morphism $\pi: X \to Y$, we obtain an isomorphism
\begin{equation}\label{decomp}
\source{hitchin-fibrations-abelian-surfaces-pw}
R\pi_*\BQ_X[\dim X- { R}]\simeq \bigoplus_{i=0}^{2{ R}}\mathcal{P}_i[-i] \in D^b_c(Y)
\end{equation}
with $\mathcal{P}_i  \in \mathrm{Perv}(Y)$. The perverse filtration can be identified as
\[
P_kH^d(X,\BQ)=\mathrm{Im}\Big\{H^{d-(\dim X - { R})}(Y, \bigoplus_{i=0}^k\mathcal{P}_i[-i])\to H^d(X,\BQ)\Big\}.
\]
In general, the isomorphism (\ref{decomp}) in the decomposition theorem is not canonical. Once we fix such an isomorphism $\phi$, we obtain a splitting $G_\ast H^\ast(X, \BQ)$ of the perverse filtration,
\[
P_kH^d(X, \BQ) = \bigoplus_{i\leq k} G_iH^d(X, \BQ).
\]
Here
\[
G_iH^d(X, \BQ) = \mathrm{Im}\Big\{ H^{d-(\dim X - { R})}(Y, \mathcal{P}_i[-i])\to H^d(X,\BQ)\Big\}
\]
with
\[
H^{d-(\dim X - {R})}(Y, \mathcal{P}_i[-i])\to H^d(X,\BQ)
\]
the morphism induced by $\phi^{-1}$.

\subsection{Perverse filtration for projective bases}
We review another description of the perverse filtration associated with
\[
\pi: X \to Y,
\]
when $Y$ is projective.

We fix $\eta$ to be an ample class on $Y$, and we consider
\[
L = \pi^* \eta \in H^2(X, \BQ).
\]
The class $L$ acts on $H^*(X, \BQ)$ as an nilpotent operator via cup product
\[
L: H^*(X, \BQ) \xrightarrow{\cup L} H^*(X, \BQ).
\]

The following proposition shows that the filtration (\ref{Perv_Filtration}) is completely described by an ample class on the base.  It is stated in the special case needed in this paper.


\begin{prop}[\cite{dCM0} Proposition 5.2.4.]
\source{hitchin-fibrations-abelian-surfaces-pw}
\notready \label{prop1.1}
\source{hitchin-fibrations-abelian-surfaces-pw}
Assume that
\begin{equation*}
\dim X = 2 \dim Y = 2R.
\end{equation*}
Then we have
\[
P_kH^m(X, \BQ) = \sum_{i\geq 1} \left(
\mathrm{Ker}(L^{R+k+i-m}) \cap \mathrm{Im}(L^{i-1}) \right) \cap H^m(X, \BQ).
\]
%Here the summation in the right-hand side stands for the span of all the vector spaces in its terms.
\end{prop}



\subsection{Abelian surfaces and Hilbert schemes}\label{Section1.5}
\source{hitchin-fibrations-abelian-surfaces-pw}
Let $A = E \times E'$ be an abelian surface with $E$ and $E'$ elliptic curves, and let $n \in \BZ^{\geq 0}$.
The natural projection, 
\[
p: A \to E'
\]
induces a morphism
\[
p_n: A^{[n]} \to {E'}^{(n)}
\]
from the Hilbert scheme of $n$ points on $A$ to the symmetric product of the elliptic curve $E'$. 
We briefly review the construction \cite{SZ} of a  \emph{canonical} splitting 
\begin{equation}\label{canonical_split}
\source{hitchin-fibrations-abelian-surfaces-pw}
H^\ast(A^{[n]}, \BQ) = \bigoplus_{i,d}\widetilde{G}_iH^d(A^{[n]},\BQ)
\end{equation}
of the perverse filtration associated with $p_n$, which culminates with  the explicit formula (\ref{cansplitting15}).

We start with a canonical decomposition 
\begin{equation}\label{split1}
\source{hitchin-fibrations-abelian-surfaces-pw}
    H^*(A, \BQ) = \bigoplus_{i,d} \widetilde{G}_iH^d(A, \BQ)
\end{equation}
with $\widetilde{G}_iH^d(A, \BQ)$ the K\"unneth factor 
\[
H^i(E, \BQ) \otimes H^{d-i}(E',\BQ) \subset H^d(A, \BQ).
\]
Since $p: A\to E'$ is a trivial fibration, (\ref{split1}) splits the perverse filtration associated with $p$. Assume that we already have decompositions of $H^*(X_1, \BQ)$ and $H^*(X_2, \BQ)$, a decomposition of $H^\ast(X_1 \times X_2, \BQ)$ can be constructed by using the K\"unneth decomposition. In particular, we obtain the direct sum decomposition of $H^*(A^n, \BQ)$ with summands
\[
\widetilde{G}_kH^*(A^n, \BQ)= \left \langle \alpha_1 \boxtimes \cdots \boxtimes \alpha_n;~~ \alpha \in \widetilde{G}_{k_i}H^*(A, \BQ),~~ \sum_{i} {k_i} = k   \right \rangle,
\]
whose $\mathfrak{S}_n$-invariant part induces a decomposition 
\[
H^*(A^{(n)}, \BQ) = \bigoplus \widetilde{G}_iH^d(A^{(n)}, \BQ).
\]
In turn, by using  K\"unneth decompositions again, this gives us  canonical decompositions for
\[
H^\ast(A^{(n_1)}\times A^{(n_2)} \times \cdots \times A^{(n_k)}, \BQ), \quad n_i \geq 1.
\]
Finally, the cohomology of the Hilbert scheme $A^{[n]}$ is related to the cohomology of symmetric products by \cite{Go2, dCM3}.  We employ exponential notation for partitions $\nu:= (\nu_i)_{i=1}^l$, with  $\nu_i >0$ and
$\sum_i \nu_i =n,$ of a positive integer $n$. Namely, we let $a_k$ be the number of times the integer $k$ appears
in the partition, so that the expression 
\[
1^{a_1}2^{a_2}\cdots n^{a_n}\]
 recovers the partition $\nu$.
We use $A^{(\nu)}$ to denote the variety $A^{(a_1)} \times A^{(a_2)} \times \cdots \times A^{(a_n)}$. The cohomology group $H^d(A^{[n]}, \BQ)$ admits a canonical decomposition
\begin{equation}\label{GS_decomp}
\source{hitchin-fibrations-abelian-surfaces-pw}
H^d(A^{[n]}, \BQ) = \bigoplus_{\nu} H^{d+2l(\nu)-2n}(A^{(\nu)}, \BQ)
\end{equation}
where $\nu$ runs through all partitions of $n$ and $l(\nu)$ is the length of $\nu$. 
The desired canonical splitting (\ref{canonical_split}) is then defined by setting
\begin{equation}\label{cansplitting15}
\source{hitchin-fibrations-abelian-surfaces-pw}
\widetilde{G}_kH^d(A^{[n]}, \BQ) = \bigoplus_{\nu} \widetilde{G}_{k+l(\nu)-n}H^{d+2l(\nu)-2n}(A^{(\nu)}, \BQ)
\end{equation}
to be the sub-vector space of $H^d(A^{[n]}, \BQ)$ under the identification (\ref{GS_decomp}). By \cite[Proposition 4.12]{Z}, this decomposition splits the perverse filtration associated with $p_n$.\footnote{We will give another interpretation of the splitting $\tilde{G}_*H^*(A^{[n]},\BQ)$ in Corollary \ref{cor3.6}.}

 In view of \cite[Proposition 2.1]{Z},  by using the K\"unneth decomposition again, we obtain a canonical splitting 
\begin{equation}\label{splitting0}
\source{hitchin-fibrations-abelian-surfaces-pw}
H^\ast(A\times A^{[n]}, \BQ) = \bigoplus_{i,d} \widetilde{G}_i H^d(A \times A^{[n]}, \BQ)
\end{equation}
of the perverse filtration associated with
\[
p \times p_n: A \times A^{[n]} \to E' \times E'^{(n)}.
\]


The following theorem obtained in \cite{Z, SZ} is concerned with  tautological classes and  multiplicativity in the context of Hilbert schemes. 

\begin{thm}
\source{hitchin-fibrations-abelian-surfaces-pw}
\notready\label{taut_Hilb}
\source{hitchin-fibrations-abelian-surfaces-pw}
Let $\CI_n$ be the universal ideal sheaf on $A \times A^{[n]}$. 
\begin{enumerate}
    \item[(a)](Tautological classes) We have
    \[
    \mathrm{ch}_k(\CI_n) \in \widetilde{G}_k H^{2k}(A\times A^{[n]}, \BQ).
    \]
    \item[(b)](Multiplicativity) The decomposition  (\ref{canonical_split}) is multiplicative, i.e.
    \[
   \cup:  \widetilde{G}_i H^d(A^{[n]}, \BQ) \times \widetilde{G}_{i'} H^{d'}(A^{[n]}, \BQ) \to \widetilde{G}_{i+i'}H^{d+d'}(A^{[n]}, \BQ).
    \]
\end{enumerate}
\end{thm}

\begin{proof}
Theorem \ref{taut_Hilb} (b) was proven in \cite{Z}. Although the multiplicativity was shown only for the perverse filtration in \cite{Z}, the same proof works for  the decomposition (\ref{canonical_split}).  This was explained in \cite[Proposition 1.8]{SZ}. Theorem \ref{taut_Hilb} (a) follows directly from the proof of \cite[Theorem 0.1]{SZ} and the proof of \cite[Theorem 3.3]{SZ}, where it is explained how to treat perverse decompositions instead of perverse filtrations. 
\end{proof}

%\subsection{Comparisons} 
%Later...
%The second cohomology of the Hilbert scheme $A^{[n]}$ admits the decomposition
%\begin{equation}\label{H2_Hilb}
%H^2(A^{[n]}, \BQ) = H^2(A, \BQ) \oplus \langle e \rangle.
%\end{equation}
%Here $e$ is the class of the exceptional divisor of the Hilbert-Chow morphism $A^{[n]} \to A^{(n)}$.

%We will show that the canonical splitting (\ref{canonical_split}) coincides with the splittings using some relative ample..
